\documentclass[12pt]{article}
\usepackage[utf8]{inputenc}
\usepackage[T1]{fontenc}
\usepackage[margin=1in]{geometry}
\usepackage{setspace}
\usepackage{xcolor}
\usepackage[colorlinks=true, urlcolor=blue!60!black, linkcolor=blue!60!black]{hyperref}
\setlength{\parindent}{2em}
\setlength{\parskip}{1.25em}
\renewcommand{\baselinestretch}{1.0}
\title{\textit{Chaplinsky v. New Hampshire}, 315 U.S. 568 (1942)}
\author{Case Brief}
\date{}
\begin{document}
\maketitle

\section{Case Metadata}
\begin{description}
  \item[Case Name \& Citation:] \textit{Chaplinsky v. New Hampshire}, 315~U.S.~568 (1942) (No.~255).
  \item[Argued / Decided:] Argued February~5, 1942; Decided March~9, 1942.
  \item[Vote:] 9--0 (affirming \textit{State v. Chaplinsky}, 91~N.H.~310, 18~A.2d~754 (1941)).
  \item[Author:] Justice Murphy (unanimous opinion joined by Chief Justice Stone and Justices Roberts, Black, Reed, Frankfurter, Douglas, Byrnes, and Jackson).
  \item[Concurrences / Dissents:] None.
  \item[Chapter / Topic:] Chapter~8: Freedom of Speech --- Categorical Exceptions to the First Amendment; Fighting Words; Low-Value Speech.
  \item[Primary Source URLs:] \href{https://supreme.justia.com/cases/federal/us/315/568/}{Justia U.S. Supreme Court Center}; \href{https://www.law.cornell.edu/constitution-conan/amendment-1/fighting-words}{Cornell LII Constitution Annotated (Fighting Words)}.
\end{description}

\section{Facts and Procedural History}
On a busy Saturday afternoon in Rochester, New Hampshire, Walter Chaplinsky, a member of the Jehovah's Witnesses, distributed religious literature on a public sidewalk and denounced organized religion as a ``racket.'' Local citizens complained to City Marshal Bowering, who informed them that Chaplinsky was lawfully engaged but warned Chaplinsky that the crowd was growing restless. After a disturbance broke out, a traffic officer escorted Chaplinsky toward the police station without formally arresting him. En route, on the public sidewalk of Wakefield Street near City Hall, they encountered Marshal Bowering hurrying toward the reported riot. After Bowering repeated his earlier warning (and, according to Chaplinsky's testimony, cursed him when Chaplinsky demanded the arrest of the crowd), Chaplinsky called Bowering ``a God damned racketeer'' and ``a damned Fascist,'' adding that ``the whole government of Rochester are Fascists or agents of Fascists.'' Chaplinsky later admitted speaking the words charged except for invoking the name of God.

Chaplinsky was charged and convicted in the Rochester municipal court under Chapter~378, \S~2 of the Public Laws of New Hampshire, which provided: ``No person shall address any offensive, derisive or annoying word to any other person who is lawfully in any street or other public place, nor call him by any offensive or derisive name\ldots'' On appeal, a jury in the New Hampshire Superior Court convicted him following a trial \textit{de novo}; the trial court excluded as immaterial evidence regarding Chaplinsky's religious mission, mistreatment by the crowd, and alleged police neglect. The Supreme Court of New Hampshire affirmed, 91~N.H.~310, 18~A.2d~754 (1941), construing the statute narrowly to reach only face-to-face ``fighting words'' likely to provoke an average addressee to immediate violence. Chaplinsky appealed to the U.S. Supreme Court, contending that the statute violated his Fourteenth Amendment freedoms of speech, press, and worship and was unconstitutionally vague.

\section{Issue(s) Presented}
Whether a state statute, authoritatively construed to forbid addressing face-to-face ``fighting words'' likely to provoke an average person to retaliation in a public place, violates the Free Speech Clause of the First Amendment (as incorporated through the Fourteenth Amendment) or is void for vagueness under the Due Process Clause.

\section{Holding}
No. Certain ``well defined and narrowly limited classes of speech''---including ``fighting words'' that ``by their very utterance, inflict injury or tend to incite an immediate breach of the peace''---are of such slight social value as a step to truth that any expressive benefit is clearly outweighed by the social interest in order and morality. As narrowly construed by the New Hampshire Supreme Court and applied to Chaplinsky's epithets (``damned racketeer'' and ``damned Fascist''), the statute is neither unconstitutionally vague nor an impermissible abridgment of free speech.

\section{Reasoning}
Writing for a unanimous Court, Justice Murphy first clarified that only free speech was genuinely implicated: the case involved spoken words rather than the press, and ``cursing a public officer'' is not an exercise of religion. Turning to the merits, the Court established a landmark two-tier categorical framework:
\begin{enumerate}
  \item \textbf{Unprotected Categories of Speech:} Free speech is ``not absolute at all times and under all circumstances.'' Drawing on Zechariah Chafee's scholarship, Murphy identified specific, narrowly limited classes of speech whose punishment had ``never been thought to raise any Constitutional problem'': the lewd and obscene, the profane, the libelous, and insulting or ``fighting'' words. Because such utterances form ``no essential part of any exposition of ideas,'' the social interest in order and morality outweighs any minimal expressive value.
  \item \textbf{Binding State Narrowing Construction:} Accepting the New Hampshire Supreme Court's severability holding and authoritative limiting construction, the Court evaluated only the statute's first clause as narrowed to face-to-face words having ``a direct tendency to cause acts of violence by the persons to whom, individually, the remark is addressed.'' The state court's test was objective: what ``men of common intelligence'' would recognize as words likely to cause an average addressee to fight.
  \item \textbf{Rejection of Vagueness and As-Applied Challenges:} Because the narrowed statute defined and punished specific ``verbal acts'' directly tending to breach the peace, it provided fair notice and satisfied due process. Taking judicial notice that calling someone a ``damned racketeer'' and ``damned Fascist'' on a public street in wartime 1942 was likely to provoke retaliation, the Court upheld the conviction and held that the Constitution did not require admitting provocation or truth as defenses.
\end{enumerate}

\section{Concurring and Dissenting Opinions}
None. \textit{Chaplinsky} was a unanimous 9--0 decision authored by Justice Murphy, reflecting broad wartime consensus on the Court that certain narrow categories of abusive or order-threatening speech lay outside the First Amendment's core.

\section{Key Quotations}
\begin{quote}
``Allowing the broadest scope to the language and purpose of the Fourteenth Amendment, it is well understood that the right of free speech is not absolute at all times and under all circumstances. There are certain well defined and narrowly limited classes of speech, the prevention and punishment of which have never been thought to raise any Constitutional problem.'' \\
--- Justice Murphy, \href{https://supreme.justia.com/cases/federal/us/315/568/}{315~U.S.~at~571--572}.
\end{quote}

\begin{quote}
``These include the lewd and obscene, the profane, the libelous, and the insulting or `fighting' words---those which, by their very utterance, inflict injury or tend to incite an immediate breach of the peace. It has been well observed that such utterances are no essential part of any exposition of ideas, and are of such slight social value as a step to truth that any benefit that may be derived from them is clearly outweighed by the social interest in order and morality.'' \\
--- Justice Murphy, \href{https://supreme.justia.com/cases/federal/us/315/568/}{315~U.S.~at~572}.
\end{quote}

\begin{quote}
``The test is what men of common intelligence would understand would be words likely to cause an average addressee to fight.'' \\
--- Justice Murphy (quoting the New Hampshire Supreme Court), \href{https://supreme.justia.com/cases/federal/us/315/568/}{315~U.S.~at~573}.
\end{quote}

\begin{quote}
``Argument is unnecessary to demonstrate that the appellations `damned racketeer' and `damned Fascist' are epithets likely to provoke the average person to retaliation, and thereby cause a breach of the peace.'' \\
--- Justice Murphy, \href{https://supreme.justia.com/cases/federal/us/315/568/}{315~U.S.~at~574}.
\end{quote}

\section{Doctrinal Significance and Modern Legacy (Through 2026)}
\textit{Chaplinsky} is the foundational source of the Supreme Court's categorical approach to the First Amendment, yet its two legacies have diverged sharply over eight decades: the ``fighting words'' doctrine itself has been relentlessly narrowed (the Supreme Court has never affirmed a fighting-words conviction since 1942), whereas \textit{Chaplinsky}'s categorical architecture---recast as a strictly historical inquiry---remains central to Roberts Court jurisprudence through 2026.
\begin{itemize}
  \item \textbf{Erosion of the Fighting Words Doctrine (\textit{Cohen} and \textit{Gooding}):} In \href{https://supreme.justia.com/cases/federal/us/403/15/}{\textit{Cohen v. California}, 403~U.S.~15 (1971)}, the Court reversed a conviction for wearing a ``Fuck the Draft'' jacket in a courthouse, limiting fighting words to ``personally abusive epithets'' directed face-to-face at a specific individual rather than general profanity. One year later, in \href{https://supreme.justia.com/cases/federal/us/405/518/}{\textit{Gooding v. Wilson}, 405~U.S.~518 (1972)}, the Court struck down Georgia's ``opprobrious words or abusive language'' statute as facially overbroad because state courts had not restricted it strictly to words tending to incite an \textit{immediate} breach of the peace---effectively discarding \textit{Chaplinsky}'s first prong (``by their very utterance inflict injury''). Justice Powell's concurrence in \href{https://supreme.justia.com/cases/federal/us/415/130/}{\textit{Lewis v. City of New Orleans}, 415~U.S.~130 (1974)}, further suggested that police officers are trained to exercise higher restraint than ordinary citizens, casting doubt on whether Chaplinsky's own conviction would survive today.
  \item \textbf{No Viewpoint Discrimination Within Fighting Words (\textit{R.A.V.}):} In \href{https://supreme.justia.com/cases/federal/us/505/377/}{\textit{R.A.V. v. City of St. Paul}, 505~U.S.~377 (1992)}, the Court held that proscribable categories under \textit{Chaplinsky} are not ``entirely invisible to the Constitution'': government may not selectively ban only race- or religion-based fighting words while permitting other fighting words.
  \item \textbf{From Value Balancing to Historical Categories (\textit{Stevens}):} In \href{https://supreme.justia.com/cases/federal/us/559/460/}{\textit{United States v. Stevens}, 559~U.S.~460 (2010)}, Chief Justice Roberts rejected the government's effort to use \textit{Chaplinsky}'s cost-benefit language to create a new unprotected category for depictions of animal cruelty. Calling ad hoc value-balancing ``startling and dangerous,'' \textit{Stevens} reinterpreted \textit{Chaplinsky} as merely \textit{describing} a closed set of ``historic and traditional categories long familiar to the bar.''
  \item \textbf{Mens Rea and Historical-Categories Enforcement (2023--2026):} In \href{https://supreme.justia.com/cases/federal/us/600/22-138/}{\textit{Counterman v. Colorado}, 600~U.S.~66 (2023)}, the Court reaffirmed the \textit{Chaplinsky}/\textit{Stevens} list of historic exceptions (including fighting words and true threats) while holding that even within unprotected true threats, the First Amendment requires a subjective recklessness \textit{mens rea} to avoid chilling protected speech. Most recently, in \href{https://supreme.justia.com/cases/federal/us/607/24-539/}{\textit{Chiles v. Salazar}, 607~U.S.~\_\_\_ (Mar.~31, 2026)}, the Court again invoked the narrow \textit{Chaplinsky}/\textit{Stevens} canon (fraud, defamation, and fighting words) to reject a state's claim that professional counseling speech constitutes a historically unprotected category.
\end{itemize}

\section{Pedagogical Critiques and Exam Hypotheticals}
\textbf{Critiques:} Modern scholars identify three major flaws in \textit{Chaplinsky}. First, its facts are notoriously ironic: Walter Chaplinsky was calling a city official a ``Fascist'' and ``racketeer'' during a heated dispute over police failure to protect an unpopular religious minority---core political criticism of government officials that modern doctrine (\textit{New York Times Co. v. Sullivan}, \textit{Cohen}, \textit{City of Houston v. Hill}) vigorously protects. Second, the Court's unexamined assumption that ``damned Fascist'' automatically triggers violence reflects an outdated, gender- and class-bound ``average addressee'' honor code that invites selective enforcement against unpopular speakers. Third, \textit{Chaplinsky}'s ``by their very utterance inflict injury'' prong collided with modern anti-censorship rules (\textit{Snyder v. Phelps}, \textit{Matal v. Tam}) holding that emotional distress alone cannot strip speech of First Amendment protection.

\textbf{Exam Hypothetical:} During a heated sidewalk protest outside a city council meeting, an activist steps within two feet of a private counter-protester, clenches his fists, and screams a string of face-to-face personal profanities and a challenge to ``settle this right here in the street.'' The activist is arrested under a state statute---authoritatively construed by the state supreme court to reach only ``face-to-face words directly tending to provoke an immediate violent breach of the peace by the person addressed.'' He challenges his conviction under the First Amendment.
\begin{itemize}
  \item \textbf{State's Argument (Under \textit{Chaplinsky}):} Unlike the general political jacket in \textit{Cohen} or the un-narrowed statute in \textit{Gooding}, this statute carries an authoritative narrowing construction matching \textit{Chaplinsky}'s breach-of-the-peace prong, is content-neutral across all topics (avoiding \textit{R.A.V.}), and was applied to a direct, face-to-face personal provocation directed at a private citizen---not a trained police officer---in a manner likely to trigger immediate physical retaliation.
  \item \textbf{Opposite Side Contrast (Distinguishing \textit{Chaplinsky}):} If the activist instead shouted ``Fascist pigs!'' from across the street through a bullhorn at uniformed police officers---or posted the identical personal insults on social media---the conviction would fail under modern post-\textit{Chaplinsky} law. Without a direct, face-to-face personal confrontation creating an imminent risk of retaliation (\textit{Cohen}; \textit{Gooding}), and especially when directed at police officers expected to exercise restraint (\textit{Lewis}) or online where immediate physical brawls cannot occur, the speech falls outside fighting words and can be punished only if it satisfies the stringent requirements of \textit{Brandenburg} incitement or \textit{Counterman} true threats (requiring subjective recklessness).
\end{itemize}

\end{document}
